\itemhead{p.47 练习 · 第3题（教材p.47）}{XJ-B1-C02-PR-P047-Y1435}
\label{XJ-B1-C02-PR-P047-Y1435}
{\footnotesize 湘教版教材p.47 / PDF第54页 / p.47 练习 / 第3题}
\field{原题}{已知二次函数 $y=x^2-2mx+m^2+3$（$m$ 为常数），求证：不论 $m$ 为何值，该二次函数均没有零点。}
\field{参考答案}{$y=(x-m)^2+3\ge3>0$，故无零点；或算得 $\Delta=-12<0$。}
{\small \field{关联知识}{Q01 二次函数、零点、判别式与根系关系（教材帮p.65 / PDF74）}
\field{方法与易错点}{按已知零点、顶点或系数选择交点式、顶点式或一般式；根平方和用根系关系消元。 易错：遗漏 $a
e0$；把二重根算成两个不同零点；有两正根只查积正。}
\field{答案依据}{答案册相关题组：PDF第13页（册内页12）；按当前题库原题逐项重算，不按旧版题号机械对套。}}
